\documentclass[12pt]{article}
\usepackage[utf8]{inputenc}
\usepackage[T1]{fontenc}
\usepackage[margin=1in]{geometry}
\usepackage{setspace}
\usepackage{xcolor}
\usepackage[colorlinks=true, urlcolor=blue!60!black, linkcolor=blue!60!black]{hyperref}
\setlength{\parindent}{0em}
\setlength{\parskip}{0.6em}
\renewcommand{\baselinestretch}{1.0}
\begin{document}
\begin{center}
{\LARGE \textit{Medell\'{i}n v. Texas}, 552 U.S. 491 (2008)\par}
\vspace{0.5em}
{\large Case Brief --- Two-Page Edition\par}
\end{center}

\section{Memory Jogger}
Neither the ICJ's \textit{Avena} judgment nor a presidential memorandum was binding domestic law; non-self-executing treaty obligations require congressional implementation (\href{https://supreme.justia.com/cases/federal/us/552/491/}{552 U.S.\ 491}).

\section{Detailed Case Facts}
Jos\'{e} Ernesto Medell\'{i}n, a Mexican national, confessed to the 1993 rape and murder of two Houston teenagers without being told of his Vienna Convention right to contact his consulate. The ICJ's \textit{Avena} judgment required the United States to provide ``review and reconsideration'' of his conviction, and President Bush directed state courts to give it effect.

\section{Procedural History}
Texas courts held his Convention claim procedurally defaulted, and the Texas Court of Criminal Appeals dismissed his second habeas application despite \textit{Avena} and the memorandum.

\section{Judicial Votes}
Affirmed, 6--3. Roberts wrote for Scalia, Kennedy, Thomas, and Alito; Stevens concurred in the judgment; Breyer, joined by Souter and Ginsburg, dissented.

\section{Holding}
\textit{Avena} binds the United States internationally, but neither it nor Article 94 of the U.N.\ Charter is self-executing domestic law. The President, acting in \textit{Youngstown}'s Category Three, cannot unilaterally make a non-self-executing treaty preempt state law.

\section{Analysis of Opinions}
Roberts read ``undertakes to comply'' as a commitment for political action, noted the Charter's diplomatic Security Council remedy, and relied on signatories' practice. The President's interests were ``plainly compelling,'' but converting a treaty into domestic law conflicted with the Senate's understanding and belongs to Congress; claims-settlement cases like \textit{Dames \& Moore} did not apply. Stevens agreed on the result but urged Texas to comply voluntarily. Breyer proposed a practical, context-specific test for self-execution and warned that some 70 treaties were jeopardized.

\section{Examples: Future Applications}
\textbf{Same-Side:} a presidential directive enforcing a non-self-executing tribunal ruling against California law fails. \textbf{Opposite-Side:} implementing legislation, like the ICSID statute, makes awards binding. \textbf{Fence-Sitter:} a \S\ 1983 damages suit for denial of consular notice.

\section{Critique}
\textit{Medell\'{i}n} lets one state put the Nation in breach of a treaty absent congressional action, and critics say it creates a de facto presumption against self-execution.

\section{Key Quotations}
\begin{quote}
``In sum, while treaties `may comprise international commitments . . . they are not domestic law unless Congress has either enacted implementing statutes or the treaty itself conveys an intention that it be ``self-executing'' and is ratified on these terms.' ''
\end{quote}
\noindent --- Chief Justice Roberts, 552 U.S.\ at 505.

\end{document}
